% Rehearsal script synchronized to conference-30min.tex.
% Check:
% python3 ../../scripts/check_script.py script.tex \
%   --deck conference-30min.tex --slot-minutes 30
\documentclass[11pt]{article}
\usepackage[margin=0.78in,headheight=14pt]{geometry}
\usepackage{xcolor}
\usepackage{microtype}
\usepackage{needspace}
\usepackage{setspace}
\usepackage{fancyhdr}
% Shared title metadata and source-traceable result phrases for the deck/script.
\newcommand{\PaperTitle}{Broadband Expansion and Small-Firm Productivity}
\newcommand{\PaperSubtitle}{Evidence from a Staggered Municipal Rollout}
\newcommand{\PaperFullTitle}{\PaperTitle: \PaperSubtitle}
\newcommand{\PaperShortTitle}{Broadband and Small-Firm Productivity}
\newcommand{\PaperAuthors}{Astrid Lindqvist \and Tomas Berg}
\newcommand{\ScriptAuthors}{Astrid Lindqvist and Tomas Berg}
\newcommand{\AuthorShort}{Lindqvist and Berg}
\newcommand{\VenueDate}{July 2026}
\newcommand{\TitleDisclaimer}{Demonstration manuscript; fictional setting and simulated data.}

\newcommand{\ResMunicipalities}{284} % Program description, paper pp. 1 and 4.
\newcommand{\ResFirms}{9{,}400} % Sample construction, paper pp. 1 and 5.
\newcommand{\ResFirmYears}{112{,}800} % Table 1 and sample construction, paper p. 5.
\newcommand{\ResYears}{2008--2019} % Sample construction, paper pp. 1 and 5.
\newcommand{\ResRolloutYears}{2011--2018} % Program description, paper pp. 1 and 4.

\newcommand{\ResBaseline}{0.034} % Table 2, col. (1), paper p. 7; log points.
\newcommand{\ResBaselineSE}{0.002} % Table 2, col. (1), paper p. 7.
\newcommand{\ResControls}{0.033} % Table 2, col. (2), paper p. 7; log points.
\newcommand{\ResControlsSE}{0.002} % Table 2, col. (2), paper p. 7.
\newcommand{\ResTrends}{0.030} % Table 2, col. (3), paper p. 7.
\newcommand{\ResTrendsSE}{0.002} % Table 2, col. (3), paper p. 7.
\newcommand{\ResMunicipalCluster}{0.034} % Table 2, col. (4), paper p. 7.
\newcommand{\ResMunicipalClusterSE}{0.003} % Table 2, col. (4), paper p. 7.

\newcommand{\ResInteractionIntercept}{0.020} % Table 4, col. (1), paper p. 9.
\newcommand{\ResInteractionInterceptSE}{0.003} % Table 4, col. (1), paper p. 9.
\newcommand{\ResInteraction}{0.032} % Table 4, col. (1), paper p. 9.
\newcommand{\ResInteractionSE}{0.005} % Table 4, col. (1), paper p. 9.
\newcommand{\ResCommMean}{0.42} % Table 1, paper p. 5.
\newcommand{\ResCommSD}{0.20} % Table 1, paper p. 5.
\newcommand{\ResCommSDGap}{0.0064} % Derived: 0.032 x 0.20; Tables 1 and 4, pp. 5 and 9.

\newcommand{\ResPlacebo}{-0.004} % Table 3, paper p. 8; log points.
\newcommand{\ResPlaceboSE}{0.002} % Table 3, paper p. 8.
\newcommand{\ResDropEarly}{0.037} % Table 5, col. (3), paper p. 9.
\newcommand{\ResDropEarlySE}{0.002} % Table 5, col. (3), paper p. 9.
\newcommand{\ResPolicyValue}{410 million euros} % Section 7, paper pp. 8--9.
\newcommand{\ResPolicyCost}{2.3 billion euros} % Section 7, paper p. 9.

\newcommand{\ResultHeadlineSlide}{Labor productivity is about 3.4\% higher after backbone availability.}
\newcommand{\ResultHeadlineSpoken}{the same firm produces about 3.4 percent more real value added per worker}
\newcommand{\ResultSupportSlide}{A one-SD increase in communication intensity adds \ResCommSDGap\ log points, about one-fifth of the main difference.}
\newcommand{\ResultSupportSpoken}{a one-standard-deviation increase in communication intensity adds about 0.0064 log points, or one-fifth of the main difference}
\newcommand{\ResultImplicationSlide}{digital infrastructure may matter most where communication frictions bind.}
\newcommand{\ResultImplicationSpoken}{digital infrastructure may matter most where communication frictions bind}


\definecolor{ScriptBlue}{HTML}{1F4E79}
\definecolor{ScriptGray}{HTML}{666666}
\setstretch{1.05}
\setlength{\parindent}{0pt}
\setlength{\parskip}{0.45em}
\widowpenalty=10000
\clubpenalty=10000
\displaywidowpenalty=10000
\pagestyle{fancy}
\fancyhf{}
\lhead{\small\textcolor{ScriptGray}{\PaperShortTitle}}
\rhead{\small\textcolor{ScriptGray}{Rehearsal script}}
\cfoot{\thepage}

\newcommand{\Pause}{\textcolor{ScriptGray}{\textit{[pause]}}}
\newcommand{\Cue}[1]{\par\nobreak\smallskip\noindent\textcolor{ScriptGray}{\footnotesize
  \textbf{Live cue:} #1}\par}
\newcommand{\scriptopening}[2]{%
  \Needspace{10\baselineskip}%
  \section*{Opening \hfill\normalfont\small\textcolor{ScriptGray}{#1}}
  #2\par}
\newcommand{\scriptframe}[3]{%
  % Keep a heading with at least the first three lines, without forcing the
  % large blank bands produced by reserving most of a page.
  \Needspace{11\baselineskip}%
  \par\medskip\noindent\textcolor{ScriptBlue}{\rule{\linewidth}{0.55pt}}%
  \par\smallskip\noindent
  \textbf{#1}\hfill\textcolor{ScriptGray}{\small #2}\par\smallskip
  #3\par}
\newcommand{\ScriptAppendix}{%
  \clearpage
  \section*{\mbox{Q\&A appendix} \hfill
    \normalfont\small\textcolor{ScriptGray}{conditional; excluded from talk time}}}

\title{Speaker script\\[0.25em]\large \PaperFullTitle}
\author{\ScriptAuthors}
\date{30-minute conference session \quad \VenueDate}

\begin{document}
\maketitle
\thispagestyle{fancy}

\scriptopening{0.5 min}{%
  Good morning. I am Astrid Lindqvist, and this is joint work with Tomas Berg.
  The setting and data are synthetic, but the economic problem is real:
  distance can keep a small firm from using outside expertise. We ask whether
  productivity is higher after broadband reaches a municipality, and whether
  the difference is greatest for firms that depend heavily on communication.
  This connects digital infrastructure to the way firms organize production.
  \Cue{Greeting; coauthor; quiet demo disclaimer; exact empirical questions.}
}

\scriptframe{Key question}{2.2 min}{%
  Think about a small manufacturer that needs an accountant, software support,
  and help coordinating with suppliers. Hiring each specialist full time makes
  little sense. A large firm can spread that fixed cost over many workers; a
  small firm usually has to buy the service outside. The problem is that
  working with a distant provider takes communication, and communication used
  to be expensive and slow.

  Broadband can change that calculation. Its value is not simply that email
  arrives faster. It can widen the set of specialists a firm can reach and make
  day-to-day coordination with a remote provider practical. That may move some
  tasks across the boundary of the firm. If a better specialist produces more
  value without a proportional increase in labor, we should see higher value
  added per worker.

  Small firms are a particularly interesting test. They have the strongest
  reason to look outside for expertise because an internal specialist is
  costly. At the same time, they may have less money and managerial capacity to
  reorganize around a new technology. So the sign and size of the productivity
  response are not obvious in advance.

  We ask two linked questions. First, is labor productivity higher after
  commercial backbone service reaches the firm's municipality? Second, is
  that difference larger for firms that began with more
  communication-intensive work? The second question is what gives the first
  one economic content. If remote coordination is the relevant friction, the
  firms with more tasks that can use it should display the larger difference.

  We observe the common date when service becomes available in each
  municipality and follow firms around that date. We do not observe each
  subscription decision. So throughout the talk, ``broadband arrival'' means
  local availability. That is the policy change we can measure consistently
  for every firm.
  \Cue{Organization problem; why small firms; two questions; availability.}
}

\scriptframe{This paper}{1.8 min}{%
  Here is what we do. We combine the rollout dates for
  \ResMunicipalities\ municipalities with annual records for \ResFirms\ small
  firms from 2008 through 2019. That gives us a long panel in which the same
  firm can be followed before and after service reaches its municipality.

  The first result is easy to interpret. In the years after local service
  arrives, \ResultHeadlineSpoken. This is a sizeable difference for a change
  in one piece of local infrastructure: the firm is producing more value from
  the same labor input. The estimate remains close to that value when we
  add firm controls, regional trends, or municipality-level clustering.

  We then build a simple task-sourcing model to ask which firms should benefit
  most. Broadband lowers the cost of coordinating with a specialist at a
  distance. A firm gains more from that change when a larger share of its work
  can use the specialist. The data line up with that prediction. A
  one-standard-deviation increase in pre-rollout communication intensity adds
  \ResCommSDGap\ log points to the post-arrival difference, about one-fifth of
  the headline association.

  The contribution is to connect a place-based infrastructure change to a
  within-firm organizational margin. The average result tells us the scale of
  the productivity pattern. The model and the cross-firm comparison tell us
  where that pattern is strongest and why. Together they suggest that
  \ResultImplicationSpoken.
  \Cue{Data link; headline result; task-sourcing prediction; organizational contribution.}
}

\scriptframe{Roadmap}{0.3 min}{%
  I will make three stops. First, the rollout and the firm data. Second, the
  timing comparison and the main productivity result. Third, the task-sourcing
  model and the cross-firm evidence that speaks to it.
  \Cue{Setting and data; design and evidence; model and interpretation.}
}

\scriptframe{Broadband reached every municipality in eight waves}{2.2 min}{%
  The rollout gives us a simple piece of timing variation. Calendar year is on
  the horizontal axis, and the vertical axis is the cumulative share of
  Nordavia's \ResMunicipalities\ municipalities with commercial backbone
  service. Every step is a new wave of municipalities coming online.

  Three dates tell most of the story. The first wave arrives in 2011, half the
  municipalities are connected by 2014, and the last wave arrives in 2018. In
  the middle of the sample, firms in one part of the country have service while
  firms elsewhere are still waiting. We use that overlap to compare firms at
  different points in the rollout. By the end, every municipality is connected,
  so the comparison comes from earlier and later arrival rather than from a
  group that never receives service.

  The 2010 Digital Infrastructure Act set the goal of universal coverage. The
  RDDF then recorded when each backbone segment entered commercial service.
  Construction followed geography---distance and terrain mattered---but also a
  policy priority score. The score moved low-productivity municipalities with
  stronger projected growth toward the front of the queue. That detail will be
  useful when we look at the event-time pattern.

  The policy object is local access to the backbone. The paper reports take-up
  above 90 percent within two years among firms with more than five employees,
  so service availability often became actual access in that group. But the
  common, consistently measured change is the municipal service date. We use
  that date as the starting point and then ask how firm outcomes move around
  it.

  The staircase is therefore more than institutional background. It creates
  the staggered timing in the empirical design. The next step is to connect
  those municipal dates to the same firms over time and to measure which firms
  had the greatest potential to use communication at a distance.
  \Cue{Axes and three dates; source of variation; selected timing; availability.}
}

\scriptframe{Data}{2.4 min}{%
  We build the analysis from three administrative sources. The business
  register tells us which firms exist and where their headquarters are. Audited
  accounts give us value added and investment. Employer--employee records give
  us employment and the occupation of each worker. After linking them, one
  observation is a firm in a year.

  The panel follows \ResFirms\ small firms from 2008 through 2019, for
  \ResFirmYears\ firm-year observations. Firms start with 3 to 49 employees and
  have positive sales. The panel is balanced, so every firm appears in all
  twelve years. This is helpful for a within-firm design: the set of firms does
  not change as we move from the pre-rollout to the post-rollout period.

  Our outcome is real value added per worker. We deflate value added using the
  producer-price index for the firm's industry, divide by full-time-equivalent
  employment, and take logs. The object is easy to interpret: how much real
  value the firm produces for each unit of labor. It captures changes in the
  productivity of the firm's overall production process, including changes in
  how inputs and tasks are organized.

  We assign each firm the commercial-service date of its headquarters
  municipality. The broadband indicator turns on in that year and stays on.
  This gives us a common measure of when the local communication constraint
  changes, even though individual subscription dates are not recorded.

  The final ingredient is communication intensity. We take each firm's 2008
  wage-bill share in occupations with above-median communication requirements.
  Client contact, supplier coordination, and remote support are examples. The
  measure is fixed in 2008, before any municipality receives service. So a firm
  that later changes its workforce cannot move itself into a different initial
  category.

  These links give us exactly the three objects the argument needs: an outcome,
  a service-arrival date, and a pre-rollout measure of the tasks that may benefit
  from cheaper remote coordination. We now use the timing to compare changes
  within firms.
  \Cue{Linked records; balanced survivor panel; outcome; exposure; 2008 task mix.}
}

\scriptframe{The design follows firms as local backbone service arrives}{2.5 min}{%
  The design asks how a firm's productivity changes when local service arrives.
  On the left of the equation is log real value added per worker. The colored
  term switches on in the service year for the firm's municipality. Beta
  compares the firm's pre- and post-arrival years with other rollout cohorts in
  the same year.

  Firm fixed effects make this a within-firm comparison. They absorb features
  such as a firm's long-run product mix, management style, and location. Year
  fixed effects absorb national shocks: a recession, a common change in input
  prices, or an economy-wide technology shift. We are therefore asking whether
  the same firm moves differently when its municipality reaches a different
  point on the rollout clock.

  The firms that are waiting provide the intuitive benchmark for firms that
  have just received service. Because all municipalities eventually connect,
  that benchmark changes over time. Early in the program there are many
  not-yet-served places; near the end there are few. The regression combines
  those timing comparisons into one summary coefficient.

  The key economic assumption is about the path without broadband. Earlier-
  and later-served municipalities should otherwise have followed similar
  productivity trends, after we remove firm and year effects. This is why the
  rollout rule matters. Places with stronger projected growth moved up the
  queue. If those projections were informative, earlier cohorts might already
  have been on a different path.

  We examine that issue directly with event time. Rather than collapse all
  post-arrival years into one indicator, the event study shows how productivity
  evolves in each year before and after service. The pre-arrival coefficients
  tell us whether the cohorts were already moving differently as arrival
  approached. The post-arrival coefficients show whether the trajectory
  changes or continues.

  Finally, service timing is assigned at the municipality level. Firms in the
  same municipality can share local shocks, so the most informative standard
  errors allow their outcomes to move together. The results table reports that
  municipality-level clustering. I will use the table for scale and the event
  study for the shape of the timing pattern.
  \Cue{Within-firm timing comparison; fixed effects; path without broadband; clustering.}
}

\scriptframe{Labor productivity is about 3.4\% higher after arrival}{2.2 min}{%
  Start with the first row. Every column uses the full panel and includes firm
  and year effects. The colored entry in the first column is \ResBaseline\ log
  points. At this magnitude, that is naturally read as about 3.4 percent higher
  real value added per worker in the post-arrival years.

  To make the comparison concrete, take the same firm before and after local
  service arrives and use firms elsewhere in the rollout to net out the common
  year. The post-arrival observations are associated with about 3.4 percent
  more output per worker. That is the economic object behind the coefficient.

  The next two columns ask how much that scale changes when we alter the
  controls. Adding firm-level covariates moves the estimate only slightly, from
  \ResBaseline\ to \ResControls. Allowing each region to have its own linear
  trend lowers it to \ResTrends. The point is not the last decimal place. The
  different ways of absorbing observed firm changes and smooth regional trends
  all leave us with a difference of roughly three percent or a little more.

  The final column changes the uncertainty calculation. It keeps the point
  estimate at \ResMunicipalCluster\ but clusters firms by municipality, which is
  where the service date varies. The standard error rises from
  \ResBaselineSE\ to \ResMunicipalClusterSE. That is the comparison I would use
  for inference because firms in the same place can share local shocks.

  So this table does two jobs. It tells us that the post-arrival productivity
  difference is economically meaningful, and it shows that its scale is not
  coming from one control set. What the table cannot show is whether the
  difference begins only when broadband arrives. For that, we need to open up
  the timing and look at the years on both sides of arrival.
  \Cue{Read the first row; 3.4 percent in levels; controls preserve scale; municipal SE; event time next.}
}

\scriptframe{Productivity rises before and after broadband arrives}{2.4 min}{%
  Now open up the timing. The horizontal axis measures years relative to the
  arrival of commercial service in a firm's municipality. Zero is the arrival
  year, and negative values are the years before it. The year just before
  arrival is the reference point. The vertical axis shows the difference in log
  labor productivity relative to that year.

  Before looking to the right, look to the left of zero. Several years before
  service, the coefficients are lower. They then climb steadily as arrival gets
  closer. The important feature is the slope, not the exact value of any point:
  productivity is already moving before the backbone becomes commercially
  available.

  After zero, the coefficients continue to rise. That is why the average
  post-arrival coefficient in the table is positive. But the post-period begins
  on a trajectory that was already underway. The clean timing story would have
  shown the earlier and later cohorts moving together before arrival and then
  separating afterward. That is not the pattern here.

  The rollout rule helps explain what we see. Municipalities with low current
  productivity but stronger projected growth received priority. If those
  projections captured real momentum, the early cohorts would be expected to
  catch up even without broadband. Their pre-arrival rise may be exactly the
  growth that placed them earlier in the queue.

  The visible pretrend limits the estimates to a descriptive reading. We see a
  clear rise in productivity around the rollout, but this timing comparison
  cannot separate the broadband response from the catch-up path used to choose
  arrival dates. That is the interpretation-changing boundary in the design,
  and this is the point in the talk where it belongs.

  There is still a useful pattern we can probe. If cheaper remote
  coordination is part of the pattern, the association should be larger for
  firms whose work relies more heavily on communication. The model turns that
  idea into a precise cross-firm prediction, and the final table asks whether
  the data follow it.
  \Cue{Orient event time; pre-period trajectory; selection logic; descriptive boundary.}
}

\scriptframe{The model predicts larger gains for communication-intensive firms}{2.5 min}{%
  The model puts one organizational choice at the center of the paper. Some
  tasks must stay inside, such as operating a machine. Accounting, software
  support, customer contact, and supplier coordination can potentially be
  bought from a remote specialist. Theta is the firm's share of tasks in that
  second group.

  A specialist can perform those tasks better; that is the quality advantage
  phi. But using the specialist has two costs. The firm pays the provider's
  price, and it pays a coordination cost, kappa, whenever information and
  decisions have to travel across distance. Broadband reduces the coordination
  cost. It makes the relationship easier to manage without changing the
  specialist's underlying price or quality.

  The inequality on the slide describes the interesting case. Before broadband,
  the quality gain is not enough to cover price plus coordination. The firm
  keeps the task inside. After broadband lowers kappa, the same outside provider
  becomes worthwhile, and the firm switches the eligible task to that
  specialist. \Pause\ The productivity gain comes from this change in who
  performs the task.

  Now compare two firms that cross the threshold. One has only a few tasks that
  can move. The other relies heavily on communication and can apply the
  specialist's quality advantage to a much larger part of production. The
  second firm gains more. That is what the positive derivative says: within the
  switching region, the productivity gain rises with theta.

  The threshold also gives us useful intuition about average effects. A firm
  that already used the specialist before broadband has no new sourcing change.
  A firm for which remote work remains too costly also does not switch. The
  response comes from firms near enough to the margin that lower coordination
  costs change the organization of production.

  In the data, theta is the firm's 2008 wage-bill share in
  communication-intensive occupations. It is fixed before rollout and captures
  how much work could plausibly benefit from coordination at a distance. The
  model's empirical prediction is simple: the post-arrival productivity
  difference should be larger when that initial share is higher.
  \Cue{Task choice; switching threshold; positive comparative static; empirical proxy.}
}

\scriptframe{The broadband association is larger for communication-intensive firms}{2.4 min}{%
  This table takes the prediction directly to the data. Both columns use log
  labor productivity and the same firm and year effects as before. The first
  column allows the post-arrival difference to vary with the firm's 2008
  communication intensity. The second gives the average post-arrival result as
  a benchmark.

  Focus on the colored interaction row. The coefficient is \ResInteraction.
  Because communication intensity is a share, a positive coefficient means
  that the productivity difference grows as we move toward firms with more
  communication-heavy work. The sign is exactly the one predicted by the task
  model.

  The broadband coefficient above it, \ResInteractionIntercept, is the fitted
  post-arrival difference for a hypothetical firm with zero communication
  intensity. For a firm elsewhere in the distribution, we add the interaction
  coefficient times its initial task share. So the economically interesting
  object is the slope across firms, not that intercept by itself.

  A one-standard-deviation comparison puts the slope in useful units. The
  standard deviation of communication intensity is \ResCommSD. Multiply that by
  \ResInteraction, and the post-arrival difference is \ResCommSDGap\ log points
  larger. In percentage terms, that is about 0.64 percentage points. Compared
  with the main 3.4 percent difference, it is roughly one-fifth as large.

  Think again about the two firms in the model. Moving one standard deviation
  toward the communication-intensive firm does not replace the average result;
  it adds about 0.64 percentage points to it. This is a comparison that lives
  well inside the observed distribution, rather than an extreme contrast
  between a firm with no communication tasks and one made up entirely of them.

  Fixing task intensity in 2008 is important for the interpretation. Firms are
  ranked by the work they did before broadband, so later workforce
  reorganization does not mechanically create the interaction. The initial
  task share delivers the cross-firm pattern the model told us to look for:
  the broadband difference is larger where communication-intensive work was
  initially more important.
  \Cue{Orient columns; positive interaction; one-SD equals 0.64 points; model prediction.}
}

\begin{samepage}
\scriptframe{Conclusion}{1.2 min}{%
  Let me close with the paper's main message. Labor productivity is about 3.4
  percent higher after broadband reaches a municipality. The difference is
  larger for firms that entered the rollout with more communication-intensive
  work: one standard deviation adds \ResCommSDGap\ log points, about one-fifth
  of the main difference.

  The model explains why that cross-firm pattern is economically meaningful.
  Better connectivity can make a distant specialist practical and move tasks
  across the boundary of the firm. Firms with more eligible tasks have more to
  gain when that coordination cost falls.

  The broader contribution is to connect public infrastructure with the
  organization of production, not only access. That is why \ResultImplicationSpoken. More generally, the value of a network
  depends on which production constraint it relaxes and which firms face that
  constraint most strongly. This broadens the policy question from coverage alone to
  the organizational frictions that determine whether firms can turn a
  connection into higher productivity.
  \Cue{3.4 percent; larger for communication-intensive firms; organization; broader implication.}
}
\end{samepage}

% A standalone single-column cue sheet for use after full-script rehearsal.
\Needspace{12\baselineskip}
\section*{Live cue layer}
\textcolor{ScriptGray}{Frame title and first spoken sentence, in deck order.}
\raggedright\small
\begin{enumerate}\setlength{\itemsep}{0.55em}\setlength{\topsep}{0.50em}
  \item \textbf{Key question} --- Think about a small manufacturer that needs an accountant, software support, and help coordinating with suppliers.
  \item \textbf{This paper} --- Here is what we do.
  \item \textbf{Roadmap} --- I will make three stops.
  \item \textbf{Broadband reached every municipality in eight waves} --- The rollout gives us a simple piece of timing variation.
  \item \textbf{Data} --- We build the analysis from three administrative sources.
  \item \textbf{The design follows firms as local backbone service arrives} --- The design asks how a firm's productivity changes when local service arrives.
  \item \textbf{Labor productivity is about 3.4\% higher after arrival} --- Start with the first row.
  \item \textbf{Productivity rises before and after broadband arrives} --- Now open up the timing.
  \item \textbf{The model predicts larger gains for communication-intensive firms} --- The model puts one organizational choice at the center of the paper.
  \item \textbf{The broadband association is larger for communication-intensive firms} --- This table takes the prediction directly to the data.
  \item \textbf{Conclusion} --- Let me close with the paper's main message.
\end{enumerate}
\normalsize

\ScriptAppendix

\scriptframe{The balanced panel describes surviving small firms}{if asked}{%
  The source requires at least four consecutive observations, then reports
  a balanced panel of \ResFirms\ firms observed in every year from 2008 through
  2019. It does not explain that added balance restriction. The final panel
  accounts exactly for \ResFirmYears\ firm-years, but it excludes
  entry, exit, and firms that fail the restrictions. Table 1 also reports an impossible
  employment combination: a mean of 23.4 and a maximum of 19. I do not use
  those employment moments in the talk. The productivity comparison therefore
  applies to selected survivors.
}

\scriptframe{Projected growth helped determine rollout order}{if asked}{%
  The 2010 Act committed the government to universal coverage by 2018. The RDDF
  then sequenced municipalities with a score that favored places below the
  median in current productivity but above the median in projected growth.
  Distance and terrain also constrained construction, but the paper does not
  isolate their independent role in timing. A natural next step would be to
  separate engineering variation from the growth score and ask whether
  productivity is stable before the resulting timing change.
}

\scriptframe{The placebo table and prose point in different directions}{if asked}{%
  Table 3 reports a placebo coefficient of \ResPlacebo\ with a standard error
  of \ResPlaceboSE. Using the rounded entries, that is about two reported
  standard errors below zero. The prose nevertheless calls it statistically
  indistinguishable from zero, and the source supplies no exact p-value or
  confidence interval. I therefore do not read the placebo as a clean
  validation or rejection of the design.
}

\scriptframe{Two numerical inconsistencies require correction}{if asked}{%
  The abstract says 34 percent, the body says 3.4 percent, and Table 2 reports
  0.034 log points. The region-trend prose also says 0.024 while Table 2 reports
  0.030. The talk therefore uses the exact table units and cells: 0.034 log
  points for the baseline and 0.030 for region trends. Both narrative values
  require correction before publication.
}

\scriptframe{Municipal assignment calls for municipal-level inference}{if asked}{%
  Treatment timing varies across \ResMunicipalities\ municipalities, while the
  baseline clusters by firm. Table 2 provides a municipality-clustered column:
  the point estimate remains \ResMunicipalCluster\ and the standard error rises
  from \ResBaselineSE\ to \ResMunicipalClusterSE. Because arrival varies by
  municipality, I would emphasize the municipality-clustered uncertainty. It
  does not solve selection in the rollout score or the visible pretrend.
}

\scriptframe{Larger gains require tasks to switch providers}{if asked}{%
  The proof covers the threshold-crossing case: remote sourcing is too expensive
  before broadband and becomes worthwhile afterward. Only then does lower
  coordination cost move eligible tasks to the specialized provider. If remote
  sourcing is worthwhile in both periods, or in neither, the sourcing decision
  does not change and this channel can generate zero gain. The proposition is
  therefore conditional.
}

\scriptframe{The interaction is consistent with, not unique to, outsourcing}{if asked}{%
  The interaction shows that the post-arrival difference is larger for firms
  with higher pre-program communication intensity. That sign matches the model.
  But the data do not observe which tasks move outside the firm, payments to
  remote providers, or task reorganization around arrival. Other characteristics
  of communication-intensive firms could generate the same pattern. Direct
  sourcing measures before and after arrival would distinguish the proposed
  mechanism from those alternatives.
}

\scriptframe{The six-year payback is not traceable from the reported inputs}{if asked}{%
  The paper compares a claimed annual flow of \ResPolicyValue\ with construction
  cost of \ResPolicyCost. But it does not report the aggregate firm value-added
  input needed to reproduce the flow. The calculation also moves from a
  labor-productivity level association to TFP, annual growth, and fiscal return
  without causality, persistence, discounting, maintenance, coverage, or
  incidence assumptions. Gross output is not government revenue. I therefore
  exclude the payback and cost-effectiveness claims from the talk's conclusion.
}

\Needspace{7\baselineskip}
\section*{Standing Q\&A stance}
Hear the whole question and count to three. Answer the strongest version in
one sentence before adding detail. Concede a genuine limitation directly,
then separate the paper's supported pattern from what remains to be shown.

\end{document}
